%------------------------------------------------------------------------------%
\documentclass[fleqn,aspectratio=169,12pt]{beamer}

\usepackage{gaal}

\title{Parábola}

%------------------------------------------------------------------------------%
\begin{document}

\addtitle

%------------------------------------------------------------------------------%
\section{Equação Canônica da Parábola}

%------------------------------------------------------------------------------%
\begin{frame}{Equação Canônica da Parábola}
  Equação \emph{canônica da parábola}
  \[
    x^2 = 4py
  \]
  com \(p \neq 0\)

  \pause
  \bigskip
  Forma alternativa
  \[
    y = \frac{1}{4p} x^2
  \]

  \pause
  \medskip
  Parábola com vértice na origem e alinhada com os eixos $x$ e $y$
\end{frame}

%------------------------------------------------------------------------------%
\begin{frame}{Equação Matricial da Parábola}
  Equação \emph{matricial da parábola}
  \[
    \pause
    v^tAv + 2B^tv = 0
  \]
  \pause
  com
  \[
    \pause
    v = \begin{bmatrix}
      x \\
      y
    \end{bmatrix}
    \qquad\qquad
    \pause
    A = \begin{bmatrix}
      1 & 0 \\
      0 & 0
    \end{bmatrix}
    \qquad\qquad
    \pause
    B = \begin{bmatrix}[c]
      0   \\
      -2p
    \end{bmatrix}
    \qquad\qquad
    \pause
    p \neq 0
  \]
\end{frame}

%------------------------------------------------------------------------------%
\begin{frame}{Elementos Geométricos da Parábola}
  \emph{Vértice} da parábola
  \[
    \mathcal{V}  = (0, 0)
  \]

  \pause
  A parábola é simétrica em relação ao eixo $y$

  \pause
  O valor de $p$ determina a abertura da parábola
  ($\abs{p}$ maior = parábola mais aberta)

  \pause
  Se \(p > 0\) a parábola abre para cima

  \pause
  Se \(p < 0 \) a parábola abre para baixo
\end{frame}

%------------------------------------------------------------------------------%
\framedgraphic*{Gráfico da Parábola}{definicao_parabola_1}{}
\framedgraphic*{Gráfico da Parábola}{definicao_parabola_2}{}
\framedgraphic*{Gráfico da Parábola}{definicao_parabola_3}{}
\framedgraphic*{Gráfico da Parábola}{definicao_parabola_4}{}
\framedgraphic*{Gráfico da Parábola}{definicao_parabola_5}{}
\framedgraphic*{Gráfico da Parábola}{definicao_parabola_6}{}
\framedgraphic*{Gráfico da Parábola}{definicao_parabola_7}{}
\framedgraphic*{Gráfico da Parábola}{definicao_parabola_8}{}

%------------------------------------------------------------------------------%
%------------------------------------------------------------------------------%
\begin{question}
  Dada a equação da parábola
  \[
    x^2 = 8y
  \]
  \vspace{-1.5\baselineskip}
  \begin{enumerate}
    \item identifique o parâmetro $p$
    \item escreva a equação matricial
    \item esboce a curva
  \end{enumerate}
\end{question}

%------------------------------------------------------------------------------%
\begin{resolution}{Parâmetro $p$}
  Comparando
  \[
    x^2 = 8y
  \]
  com a equação canônica
  \[
    x^2 = 4py
  \]
  obtemos
  \[
    4p = 8
    \qquad
    p = 2
  \]
\end{resolution}

%------------------------------------------------------------------------------%
\begin{resolution}{Equação Matricial}
  Equação Matricial
  \[
    v^tAv + 2B^tv = 0
  \]
  com
  \[
    A = \begin{bmatrix}
      1 & 0 \\
      0 & 0
    \end{bmatrix}
    \qquad\qquad
    B
    = \begin{bmatrix}[c]
      0   \\
      -2p
    \end{bmatrix}
    = \begin{bmatrix}
      0  \\
      -4
    \end{bmatrix}
  \]
\end{resolution}

%------------------------------------------------------------------------------%
\begin{resolution}{Esboço}
\end{resolution}

%------------------------------------------------------------------------------%
\endinput


%------------------------------------------------------------------------------%
\section{Parábola Simétrica no Eixo $x$}

%------------------------------------------------------------------------------%
\begin{frame}{Equações da Parábola Simétrica no Eixo $x$}
  Equação canônica da parábola \emph{simétrica no eixo $x$}
  \[
    y^2 = 4px
    \qquad\qquad
    p \neq 0
  \]
  \pause
  Forma alternativa
  \[
    x = \frac{1}{4p} y^2
  \]
  \pause
  Equação matricial:
  \(
    v^tAv + 2B^tv = 0
  \)
  \[
    \pause
    v = \begin{bmatrix}
      x \\
      y
    \end{bmatrix}
    \qquad\qquad
    \pause
    A = \begin{bmatrix}
      0 & 0 \\
      0 & 1
    \end{bmatrix}
    \qquad\qquad
    \pause
    B = \begin{bmatrix}[c]
      -2p   \\
      0
    \end{bmatrix}
  \]
\end{frame}

%------------------------------------------------------------------------------%
\framedgraphic*{Gráfico da Parábola Simétrica no Eixo $x$}{definicao_parabola_y_1}{}
\framedgraphic*{Gráfico da Parábola Simétrica no Eixo $x$}{definicao_parabola_y_2}{}
\framedgraphic*{Gráfico da Parábola Simétrica no Eixo $x$}{definicao_parabola_y_3}{}
\framedgraphic*{Gráfico da Parábola Simétrica no Eixo $x$}{definicao_parabola_y_4}{}
\framedgraphic*{Gráfico da Parábola Simétrica no Eixo $x$}{definicao_parabola_y_5}{}
\framedgraphic*{Gráfico da Parábola Simétrica no Eixo $x$}{definicao_parabola_y_6}{}
\framedgraphic*{Gráfico da Parábola Simétrica no Eixo $x$}{definicao_parabola_y_7}{}
\framedgraphic*{Gráfico da Parábola Simétrica no Eixo $x$}{definicao_parabola_y_8}{}

%------------------------------------------------------------------------------%
%------------------------------------------------------------------------------%
\begin{question}
  Dada a equação da parábola
  \[
    (x - 3)^2 = 12(y + 1)
  \]
  \vspace{-1.5\baselineskip}
  \begin{enumerate}
    \item identifique vértice e o parâmetro $p$
    \item escreva a equação matricial
    \item esboce a curva
  \end{enumerate}
\end{question}

%------------------------------------------------------------------------------%
\begin{resolution}{Parâmetro $p$}
  Comparando
  \[
    (x - 3)^2 = 12(y + 1)
  \]
  com a equação canônica transladada
  \[
    (x-h)^2 = 4p(y-k)
  \]
  obtemos
  \[
    4p = 12
    \qquad
    p = 3
  \]
  vértice:
  \(
    \mathcal{V}
    = (h, k)
    = (3, -1)
  \)
\end{resolution}

%------------------------------------------------------------------------------%
\begin{resolution}{Equação Matricial}
  Equação Matricial
  \[
    V^tAv + 2B^tV = 0
  \]
  com
  \[
    A = \begin{bmatrix}
      1 & 0 \\
      0 & 0
    \end{bmatrix}
    \qquad\qquad
    B
    = \begin{bmatrix}[c]
      0   \\
      -2p
    \end{bmatrix}
    = \begin{bmatrix}
      0  \\
      -6
    \end{bmatrix}
  \]
  \[
    V
    = \begin{bmatrix}
      x - h  \\
      y - k
    \end{bmatrix}
    = \begin{bmatrix}[c]
      x - 3  \\
      y - (-1)
    \end{bmatrix}
    = \begin{bmatrix}
      x - 3  \\
      y + 1
    \end{bmatrix}
  \]
\end{resolution}

%------------------------------------------------------------------------------%
\begin{resolution}{Esboço}
\end{resolution}

%------------------------------------------------------------------------------%
\endinput


%------------------------------------------------------------------------------%
\section{Translação da Parábola}

%------------------------------------------------------------------------------%
\begin{frame}{Translação da Parábola}
  Novo vértice \(C = (h, k)\)

  \pause
  Novas coordenadas
  \[
    \pause
    X = x - h
    \qquad
    \pause
    Y = y - k
  \]

  \pause
  Equação da parábola transladada
  \[
    \pause
    X^2 = 4pY
  \]
  \[
    \pause
    (x-h)^2 = 4p(y-k)
  \]
\end{frame}

%------------------------------------------------------------------------------%
\begin{proofframe}{Equação Matricial da Parábola Transladada}
  Novo vértice
  \[
    \pause
    C
    = \begin{bmatrix}
      h \\
      k
    \end{bmatrix}
  \]

  \bigskip
  \pause
  Novas variáveis
  \[
    \pause
    V
    \pause
    = \begin{bmatrix}
      X \\
      Y
    \end{bmatrix}
    \pause
    = v - C
    \pause
    = \begin{bmatrix}
      x \\
      y
    \end{bmatrix}
    -
    \begin{bmatrix}
      h \\
      k
    \end{bmatrix}
    \pause
    = \begin{bmatrix}
      x - h \\
      y - k
    \end{bmatrix}
  \]
\end{proofframe}

%------------------------------------------------------------------------------%
\begin{proofframe}{Equação Matricial da Parábola Transladada}
  \begin{align*}
    \onslide<+->{\emph{V^tAV + 2B^tV} & \emph{= 0} \\}
    \onslide<+->{(v-C)^tA(v-C) + 2B^t(v-C) & = 0 \\}
    \onslide<+->{\begin{bmatrix}
      x - h \\
      y - k
    \end{bmatrix}^t
    \begin{bmatrix}
      1 & 0              \\
      0 & 0
    \end{bmatrix}
    \begin{bmatrix}
      x - h \\
      y - k
    \end{bmatrix}
    +
    \begin{bmatrix}[c]
      0   \\
      -2p
    \end{bmatrix}^t
    \begin{bmatrix}
      x - h \\
      y - k
    \end{bmatrix}
    & = 0}
    \\
    \onslide<+->{\emph{(x - h)^2 -2p (y - k)} & \emph{= 0}}
  \end{align*}
\end{proofframe}

%------------------------------------------------------------------------------%
\begin{proofframe}{Equação Matricial da Parábola Transladada}
  \begin{align*}
    \onslide<+->{(x - h)^2 -2p (y - k)       & = 0 \\}
    \onslide<+->{x^2 - 2hx + h^2 - 2py + 2pk & = 0 \\}
    \onslide<+->{x^2 - 2hx - 2py + h^2 + 2pk & = 0 \\}
    \onslide<+->{\begin{bmatrix}
      x \\
      y
    \end{bmatrix}^t
    {\color<6->{magenta}\begin{bmatrix}
      1 & 0              \\
      0 & 0
    \end{bmatrix}}
    \begin{bmatrix}
      x \\
      y
    \end{bmatrix}
    +
    {\color<7->{purple} \begin{bmatrix}[c]
      -2h   \\
      -2p
    \end{bmatrix}^t}
    \begin{bmatrix}
      x \\
      y
    \end{bmatrix}
    + {\color<8->{violet}(h^2 + 2pk)}
    & = 0}
    \\
    \onslide<+->{v^t {\color<6->{magenta} A} v
    + {\color<7->{purple} H^t}v
    + {\color<8->{violet} f} & = 0}
  \end{align*}
\end{proofframe}

%------------------------------------------------------------------------------%
%------------------------------------------------------------------------------%
\begin{question}
  Determine o conjunto de soluções reais da equação
  \[
    (x-2)^2 = 0
  \]
\end{question}

%------------------------------------------------------------------------------%
\begin{resolution}{Solução}
  \begin{align*}
    (x-2)^2 & = 0 \\
    x-2     & = 0 \\
    x       & = 2
  \end{align*}
  Não há restrição para $y$

  O conjunto de soluções é a reta vertical
  \(
    x = 2
  \)
  \[
    S = \left\{ (x, y) \in \R^2 \st x = 2 \right\}
  \]
  A equação é um caso degenerado de equação quadrática no plano
\end{resolution}

%------------------------------------------------------------------------------%
\endinput

%------------------------------------------------------------------------------%
\begin{question}
  Determine e esboce o conjunto de soluções reais da equação
  \[
    (x-2)^2 = 0
  \]
\end{question}

%------------------------------------------------------------------------------%
\begin{resolution}{Solução}
  \begin{align*}
    \onslide<+->{(x-2)^2 & = 0 \\}
    \onslide<+->{x-2     & = 0 \\}
    \onslide<+->{x       & = 2}
  \end{align*}
  \onslide<+->{Não há restrição para $y$}

  \onslide<+->{O conjunto de soluções é a reta vertical
  \(
    x = 2
  \)}
  \[
    \onslide<+->{S = \left\{ (x, y) \in \R^2 \st x = 2 \right\}}
  \]
  \onslide<+->{A equação é um caso degenerado de equação quadrática no plano}
\end{resolution}

%------------------------------------------------------------------------------%
\begin{resolution}{Esboço}
  \centering
  \includegraphics{I-05-Parabola-ex-4}
\end{resolution}

%------------------------------------------------------------------------------%
\endinput


%%------------------------------------------------------------------------------%
%\section{Lista Mínima}
%
%%------------------------------------------------------------------------------%
%\begin{frame}{Lista Mínima}
%  \begin{enumerate}
%    \item
%  \end{enumerate}
%
%  \vfill
%  Atenção: A prova é baseada no livro, não nas apresentações
%\end{frame}

%------------------------------------------------------------------------------%
\end{document}
%------------------------------------------------------------------------------%
